\documentclass[11pt]{article}
\usepackage{amsmath,amssymb,amsthm}
\usepackage[T1]{fontenc}
\usepackage{amsmath,amssymb,amsthm}
\usepackage{fullpage}
\usepackage[hidelinks]{hyperref}

\newcommand{\Z}{\mathbb Z}
\newcommand{\Q}{\mathbb Q}
\newcommand{\R}{\mathbb R}
\newcommand{\F}{\mathbb F}
\newcommand{\Surg}[2]{S^3_{#1}(#2)}
\newcommand{\HFhat}{\widehat{HF}}
\newcommand{\Int}{\operatorname{int}}
\newcommand{\nbd}{\nu}
\newcommand{\rk}{d}

\newtheorem{theorem}{Theorem}[section]
\newtheorem{lemma}[theorem]{Lemma}
\newtheorem{proposition}[theorem]{Proposition}
\newtheorem{corollary}[theorem]{Corollary}
\theoremstyle{definition}
\newtheorem{definition}[theorem]{Definition}
\theoremstyle{remark}
\newtheorem{remark}[theorem]{Remark}
\newtheorem{example}[theorem]{Example}

\newcommand{\fptitle}{Review of Problem 6 -- ProofCouncil}
\newcommand{\reviewtag}{Reviewer 3}
\usepackage{fp-review}

\begin{document}
\maketitle

\section{Review}

\paragraph{Correctness.}
\begin{itemize}
	\item  1: essentially flawless mathematical argument
	\item  0.5: minor errors that are easily corrected
	\item  0: blank solution, or acknowledged partial result
	\item  $-0.5$: major errors
	\item  $-1$: proposed proof makes no progress towards the result.
\end{itemize}
Note that points should be docked for wrong statements/ solutions.

\noindent Score: 1. 

\paragraph{Exposition.}
\begin{itemize}
	\item  1: well-written and clear solution.
	\item  0: argument that is possible to follow with serious effort
	\item  $-1$: very poor exposition
\end{itemize}

\noindent Score: 1. 


\paragraph{Attributions.}
\begin{itemize}
	\item  1: solutions contains fully appropriate attributions
		to related work and theorems being used.
	\item  0: moderate issues with missing attributions or hallucinated references
	\item  $-1$: blatant plagiarism or hallucinations
\end{itemize}

\noindent Score: 0.

This is a correct and complete solution. The open case is that of hyperbolic L-space knots. The main argument there is by contradiction: assume that surgery on an L-space knot has 3 or more prime summands, then derive contradictory bounds on the absolute value of the surgery slope. The key step is to consider an appropriate covering space of the surgery and argue that any additional connected summand would make the rank of the fundamental group of this covering space too high, contradicting an upper bound that comes from the fact that the knot is fibered. The rest of the proof relies on combining appropriate results in the literature. 

The proof is reasonably well-written and quite easy to follow. There are certainly places where the exposition can be improved, for example the text refers to lemmas that will be stated in the future. Sometimes language or precision issues make it a bit hard to discern the intended meaning. However, I wouldn't say it rises to the level of  ``possible to follow with serious effort''. It would be more accurate to say the paper is annoying to read.

It's worth noting is that rather a lot of space is taken up by reviewing or reproving basic facts, e.g. knot groups have weight 1. It's by no means efficient; but, depending on the intended audience, this level of detail can be well-suited. 

For the most part, attributions are appropriate and more or less complete. There are a handful of instances where an attribution is made but a specific reference is omitted. A few times the text cites a paper referencing another paper, rather than citing the original source.  I would give it a 0.5 if this was an option.

\newpage
\section{Submission with inline comments}


\begin{abstract}
An \(L\)-space knot is a knot in \(S^3\) with a non-trivial Dehn surgery
whose Heegaard Floer homology is as small as its first homology allows.
We show that every Dehn surgery on an \(L\)-space knot has at most two
non-trivial prime summands. The proof combines published results on
reducible surgeries (\(L\)-space knots are fibered, by Ni; a reducible
surgery on a hyperbolic knot with a positive \(L\)-space surgery has slope
\(2g-1\), by Hom--Lidman--Zufelt; the Cabling Conjecture holds for
non-trivial non-hyperbolic knots, by Moser and Scharlemann; Gordon's
formula for
surgery on a cable; the Gordon--Luecke theorem on non-integral surgeries)
with an elementary group-theoretic obstruction: if a \(p/q\)-surgery on a
fibered knot of genus \(g\) has a non-trivial integer homology sphere
summand, then \(|p|\le g\). The obstruction is proved by estimating the
minimal number of generators of the commutator subgroup of the
fundamental group in two ways, through the universal abelian cover and
Grushko's theorem on free products on one side and through the
fundamental group of the fiber surface on the other.
\end{abstract}

\section{Introduction}\label{sec:intro}

Dehn surgery along a knot \(K\subset S^3\) produces a closed oriented
\(3\)-manifold \(\Surg{r}{K}\) for every slope
\(r\in\Q\cup\{\infty\}\) (notation is recalled in
Section~\ref{sec:surgery}). By the Kneser--Milnor theorem
\review{add citation here}
every closed
oriented \(3\)-manifold is a connected sum of prime manifolds, unique up to
order, and a basic question in Dehn surgery theory is how complicated this
decomposition can be for a surgery on a knot. For a non-trivial knot
\(K\subset S^3\), the Cabling Conjecture of Gonz\'alez-Acu\~na and Short
\cite{GAS1986} predicts that \(\Surg{r}{K}\) is reducible only when
\(K\) is a cable knot and \(r\) is the slope of the cabling annulus (the
unknot must be excluded, since its \(0\)-surgery is \(S^1\times S^2\)); in
that case the surgery is a lens space connected-summed
with a surgery on the companion knot (Gordon \cite{Gordon1983}), and the
latter surgery is prime, because its slope is non-integral (Gordon--Luecke
\cite{GordonLuecke1987}) or because the companion is the unknot, so the
conjecture predicts in particular that a
surgery on a knot in \(S^3\) never has three non-trivial prime summands.
The conjecture is known for
torus knots (Moser \cite{Moser1971}) and for satellite knots (Scharlemann
\cite{Scharlemann1990}), and it is open for hyperbolic knots. Without any
hypothesis on the knot, Howie \cite{Howie2002} showed that at most two
prime summands of a surgery on a knot in \(S^3\) can have non-zero first
homology, so a third summand would have to be an integer homology sphere;
see also \cite{Howie2010}.

This note answers the question for knots satisfying a strong
Heegaard Floer restriction. An \emph{\(L\)-space} is a rational homology
sphere \(Y\) whose Heegaard Floer homology \(\HFhat(Y)\) has the smallest
possible dimension, namely \(|H_1(Y;\Z)|\); lens spaces are examples. A knot
\(K\subset S^3\) is an \emph{\(L\)-space knot} if \(\Surg{r}{K}\) is an
\(L\)-space for some slope \(r\neq\infty\). We allow \(r\) of either sign;
some authors require \(r>0\), and we explain at the end of
Section~\ref{sec:Lspace} and in Step~1 of Section~\ref{sec:proof} why the
two conventions lead to the same theorem. These notions are recalled in
Section~\ref{sec:Lspace}. Torus knots are \(L\)-space knots, since
suitable surgeries on them are lens spaces \cite{Moser1971}.

We count the prime summands of a manifold other than \(S^3\) in its
Kneser--Milnor decomposition, and we say that \(S^3\) has none
(Section~\ref{sec:prime}); these are its \emph{non-trivial prime
summands}.

\begin{theorem}\label{thm:main}
Let \(K\subset S^3\) be an \(L\)-space knot. For every slope
\(r\in\Q\cup\{\infty\}\), the manifold \(\Surg{r}{K}\) has at most two
non-trivial prime summands.
\end{theorem}

The \(L\)-space hypothesis enters the proof at exactly two points, each
through a published theorem: \(L\)-space knots are fibered (Ni
\cite{Ni2007}), and a reducible surgery on a \emph{hyperbolic} knot
admitting a positive \(L\)-space surgery has slope exactly \(2g-1\), where
\(g\) is the genus of the knot (Hom--Lidman--Zufelt \cite[Theorem~1.3]{HLZ2015}).
For non-trivial non-hyperbolic knots no Floer-theoretic input is needed,
because the Cabling Conjecture is a theorem there. The argument does not identify
the reducible surgery itself, and it does not prove the Cabling
Conjecture for hyperbolic \(L\)-space knots; see Remark~\ref{rem:notproved}.
The main ingredient we prove here is the following statement about
fibered knots, established in Section~\ref{sec:rank} as
Proposition~\ref{prop:large}.

\begin{theorem}\label{thm:fibered}
Let \(K\subset S^3\) be a fibered knot of genus \(g\geq1\), and let
\(p/q\) be a slope with \(p\neq0\) and \(|p|\geq 2g-1\). Then
\(\Surg{p/q}{K}\) has at most two non-trivial prime summands.
\end{theorem}

\paragraph{Strategy.}
Suppose that \(M=\Surg{p/q}{K}\) has three non-trivial prime summands.
\review{Does this approach address the case of more than 3 prime summands? (Turns out that yes, but should say ``at least''}

By Howie's theorem one of them, call it \(\Sigma\), has
\(H_1(\Sigma;\Z)=0\), so \(M\cong X\#\Sigma\) with \(\Sigma\) an integer
homology sphere other than \(S^3\) and \(H_1(X;\Z)\cong H_1(M;\Z)\cong\Z/|p|\).
Let \(G=\pi_1(M)\) and let \(G'\) be its commutator subgroup. We compute
the minimal number of generators \(\rk(G')\) in two ways.
\begin{itemize}
\item \emph{From above}, via the fiber. The group \(G\) is a quotient of
the knot group \(\pi_1(E_K)\), where \(E_K\) is the knot exterior, whose
commutator subgroup is the
fundamental group of the fiber surface, a free group of rank \(2g\). An
epimorphism maps commutator subgroups onto commutator subgroups, so
\(\rk(G')\leq 2g\) (Lemma~\ref{lem:fiber}).
\item \emph{From below}, via the universal abelian cover. The subgroup
\(G'\) is the fundamental group of the \(|p|\)-fold cover
\(\widetilde M\to M\) with deck group \(G/G'=H_1(M;\Z)\). Because
\(\pi_1(\Sigma)\) is perfect, it maps trivially to the deck group, so the
punctured copy of \(\Sigma\) inside \(M\) lifts to \(|p|\) disjoint copies,
and \(\widetilde M\cong\widetilde X\#\Sigma^{\#|p|}\). Since
\(\pi_1(\Sigma)\) is non-trivial and perfect it is not cyclic, so it needs
at least two generators, and Grushko's theorem gives
\(\rk(G')\geq 2|p|\) (Lemma~\ref{lem:cover}).
\end{itemize}
Hence \(|p|\leq g\). This is incompatible with \(|p|\geq 2g-1\) unless
\(g=1\) and \(|p|=1\), an endpoint case that is excluded by a direct
application of Grushko's theorem to the perfect group \(G\).

\review{Attribution is present but citation is missing}

Theorem~\ref{thm:main} then follows by Thurston's trichotomy, according
to which every non-trivial knot is a torus knot, a satellite knot or
hyperbolic (Section~\ref{sec:types}). For a hyperbolic \(L\)-space knot, Ni's theorem and the Hom--Lidman--Zufelt
theorem put every reducible surgery exactly in the range of
Theorem~\ref{thm:fibered}; for a torus or satellite knot the Cabling
Conjecture, Gordon's formula for surgery on a cable, and the
Gordon--Luecke theorem that non-integral surgeries are irreducible leave
room for at most two summands.

\paragraph{Outline.}
Section~\ref{sec:background} fixes
conventions and recalls the standard facts we use,
Section~\ref{sec:rank} proves the group-rank obstruction,
Section~\ref{sec:inputs} states the published theorems that are used as
black boxes, Section~\ref{sec:proof} proves Theorem~\ref{thm:main}, and
Section~\ref{sec:comments} discusses what the argument does and does not
show.

A genus bound of the shape \(|p|\leq g\) for surgeries with a non-trivial
integer homology sphere summand is attributed to Gonz\'alez-Acu\~na and
Short \cite[Proposition~1.4]{GAS1986} in \cite[Theorem~1.2]{BodishDeYeso},
in the case where the complementary summand \(X\) is a lens space. The
case needed here has \(X\) non-prime, so we prove Lemma~\ref{lem:cover}
for an arbitrary complementary summand. Related statements about
homology sphere summands of surgeries on \(L\)-space knots appear in
\cite[Section~4]{BodishDeYeso}. We make no claim of priority for
Theorem~\ref{thm:main}; the aim of this note is a complete and
self-contained argument modulo the cited theorems.

\review{Hard to follow, e.g. refers to Lemma~\ref{lem:cover}; should instead state the relevant inequality.}

\section{Conventions and background}\label{sec:background}

\subsection{Dehn surgery}\label{sec:surgery}

Throughout, \(K\subset S^3\) is a tame knot, \(\nbd K\) is a closed
tubular neighbourhood of \(K\), and \(E_K=S^3\setminus\Int\nbd K\) is the
\emph{exterior} of \(K\), a compact oriented \(3\)-manifold with torus
boundary. On \(\partial E_K=\partial\nbd K\) we fix the \emph{meridian}
\(\mu\), a simple closed curve bounding a disc in \(\nbd K\), and the
\emph{preferred longitude} \(\lambda\), a simple closed curve isotopic to
\(K\) in \(\nbd K\) and null-homologous in \(E_K\) (it is the boundary of
a Seifert surface for \(K\)). A \emph{slope} is an isotopy class of
unoriented essential simple closed curves on \(\partial E_K\). The curve
homologous to \(p\mu+q\lambda\), with \(\gcd(p,q)=1\), determines the
slope \(p/q\in\Q\cup\{\infty\}\), where \(\infty=1/0\) is the meridional
slope; for finite slopes we normalise \(q>0\), so that
\(p\in\Z\) is determined.

The \emph{Dehn surgery} on \(K\) with slope \(r\) is the closed oriented
\(3\)-manifold
\[
 \Surg{r}{K}=E_K\cup_\varphi (S^1\times D^2),
\]
where the gluing homeomorphism
\(\varphi\colon\partial(S^1\times D^2)\to\partial E_K\) sends the meridian
\(\{*\}\times\partial D^2\) of the solid torus to a curve of slope \(r\),
and the orientation extends that of \(E_K\subset S^3\). Its
homeomorphism type depends only on \(r\). We use the following standard
facts.

\begin{enumerate}
\item[(S1)] \(\Surg{\infty}{K}=S^3\), since filling along the meridian
undoes the removal of \(\nbd K\).
\item[(S2)] \(H_1(E_K;\Z)\cong\Z\), generated by \(\mu\), and \(\lambda=0\)
in \(H_1(E_K;\Z)\). Filling along \(p\mu+q\lambda\) kills the class
\(p[\mu]\), so
\[
 H_1(\Surg{p/q}{K};\Z)\cong\Z/|p|\Z,
\]
which is to be read as \(\Z\) when \(p=0\). In particular
\(\Surg{p/q}{K}\) is a rational homology sphere if and only if \(p\neq0\).
\item[(S3)] By van Kampen's theorem,
\(\pi_1(\Surg{p/q}{K})\cong\pi_1(E_K)/\langle\!\langle\gamma_{p/q}\rangle\!\rangle\),
the quotient by the normal closure of the element \(\gamma_{p/q}\) of
\(\pi_1(E_K)\) represented by the filling curve of slope \(p/q\) (the
elements of \(\pi_1(E_K)\) represented by \(\mu\) and \(\lambda\) are
again denoted \(\mu\) and \(\lambda\)). Applied to \(p/q=\infty\) this
gives
\(\pi_1(E_K)/\langle\!\langle\mu\rangle\!\rangle\cong\pi_1(S^3)=1\): the
knot group is normally generated by a meridian.
\item[(S4)] Let \(\overline K\) be the mirror image of \(K\). A reflection
of \(S^3\) carrying \(K\) to \(\overline K\) preserves meridians and
preferred longitudes but reverses the intersection pairing on the
boundary torus, so it carries the slope \(r\) of \(K\) to the slope
\(-r\) of \(\overline K\) and induces an orientation-reversing
homeomorphism \(\Surg{r}{K}\to\Surg{-r}{\overline K}\); equivalently,
there is an orientation-preserving identification
\(\Surg{-r}{\overline K}\cong-\Surg{r}{K}\), where \(-M\) denotes \(M\)
with the opposite orientation.
\item[(S5)] For the unknot \(U\), \(\Surg{p/q}{U}\) is a lens space
with first homology of order \(|p|\) (namely \(L(p,q)\) up to
orientation) when \(|p|\geq2\), is \(S^3\) when \(|p|=1\), and is
\(S^1\times S^2\) when \(p=0\). Here, for \(p\geq2\) (negative \(p\)
differs only by orientation), \(L(p,q)\) denotes the quotient of the
unit sphere \(S^3\subset\mathbb C^2\) by the free action of \(\Z/p\Z\)
generated by \((z,w)\mapsto(\zeta z,\zeta^{q}w)\), \(\zeta=e^{2\pi i/p}\);
it has \(H_1(L(p,q);\Z)\cong\Z/p\Z\) and universal cover \(S^3\).
Orientation conventions for lens spaces play no role in this note, since
only the number of prime summands is counted.
\end{enumerate}

\subsection{Prime decomposition}\label{sec:prime}

Let \(M\) be a closed connected oriented \(3\)-manifold. \(M\) is
\emph{irreducible} if every embedded \(2\)-sphere in \(M\) bounds a
\(3\)-ball, and \emph{reducible} otherwise. \(M\) is \emph{prime} if
\(M\cong M_1\#M_2\) implies \(M_1\cong S^3\) or \(M_2\cong S^3\). An
irreducible manifold is prime, and conversely a prime closed orientable
\(3\)-manifold is either irreducible or homeomorphic to \(S^1\times S^2\)
\cite[\S1.1]{Hatcher3M}. Thus for a rational homology sphere, prime and
irreducible are the same thing, and non-prime is the same as reducible.

If a covering space of \(M\) is irreducible, so is \(M\)
\cite[Proposition~1.6]{Hatcher3M}; by Alexander's theorem
\cite[\S1.1]{Hatcher3M}, \(S^3\) and \(\R^3\) are irreducible. Hence
lens spaces, which are covered by \(S^3\), are irreducible, and so is
any closed \(3\)-manifold with universal cover \(\R^3\).

The Kneser--Milnor theorem \cite[Theorem~1.5]{Hatcher3M} states that
every closed oriented \(3\)-manifold \(M\not\cong S^3\) is a connected sum
\(M\cong M_1\#\cdots\#M_k\) of prime manifolds \(M_i\not\cong S^3\), and
that the collection \(\{M_i\}\) is unique up to order and
orientation-preserving homeomorphism. We call \(M_1,\dots,M_k\) the
\emph{non-trivial prime summands} of \(M\); by uniqueness their number
\(k\) is an invariant of \(M\), and we set \(k=0\) for \(M\cong S^3\).
Reversing the orientation of \(M\) reverses that of each \(M_i\), so
\(M\) and \(-M\) have the same number of non-trivial prime summands.
By van Kampen's theorem, for any connected sum of closed oriented
\(3\)-manifolds \(N_1,\dots,N_k\) (prime or not),
\begin{equation}\label{eq:vk}
 \pi_1(N_1\#\cdots\#N_k)\cong\pi_1(N_1)*\cdots*\pi_1(N_k),
 \qquad
 H_1(N_1\#\cdots\#N_k;\Z)\cong\bigoplus_{i=1}^k H_1(N_i;\Z),
\end{equation}
since the connected sum is formed by gluing punctured copies of the
\(N_i\) along \(2\)-spheres, and removing a ball from a \(3\)-manifold
does not change its fundamental group.

A closed connected oriented \(3\)-manifold \(\Sigma\) with
\(H_1(\Sigma;\Z)=0\) is called an \emph{integer homology sphere}; by
Poincar\'e duality and the universal coefficient theorem,
\(H_*(\Sigma;\Z)\cong H_*(S^3;\Z)\).

\subsection{Torus, satellite and hyperbolic knots; cables}\label{sec:types}

A non-trivial knot \(K\subset S^3\) is a \emph{torus knot} if it can be
isotoped onto the boundary of a standardly embedded solid torus; it is a
\emph{satellite knot} if \(E_K\) contains an incompressible torus (an
embedded torus whose fundamental group injects into \(\pi_1(E_K)\)) that
is not parallel to \(\partial E_K\); and it is \emph{hyperbolic} if
\(S^3\setminus K\) admits a complete hyperbolic metric of finite volume.
By Thurston's hyperbolization theorem, every non-trivial knot is of
exactly one of these three types \cite[Corollary~2.5]{Thurston1982}.
Composite knots are satellite knots: if \(K=K_1\#K_2\) with both
summands non-trivial, the boundary of a regular neighbourhood of the
union of \(K\) with a ball containing the \(K_2\) summand is a torus in
\(E_K\), knotted like \(K_1\), which is incompressible and not
boundary-parallel (the ``swallow--follow'' torus).

Let \(J\subset S^3\) be a knot, with meridian \(\mu_J\) and preferred
longitude \(\lambda_J\) on \(\partial\nbd J\), and let \(w\geq2\) and
\(t\) be coprime integers (so \(t\neq0\)). The \emph{\((w,t)\)-cable} of
\(J\) is the knot
\[
 C_{w,t}(J)\subset\partial\nbd J\subset S^3
\]
given by a simple closed curve on \(\partial\nbd J\) homologous to
\(t\mu_J+w\lambda_J\) there. Its image in \(H_1(\nbd J;\Z)\cong\Z\) is
\(w\) times the generator, so \(w\) is the \emph{winding number} of the
cable. If \(J\) is the unknot, \(C_{w,t}(J)\) is the torus knot
\(T_{w,t}\); if \(J\) is non-trivial, \(C_{w,t}(J)\) is a satellite knot
with companion \(J\). We call \(wt\) the \emph{cabling slope} of
\(C_{w,t}(J)\): it is the framing of the cable induced by the torus
\(\partial\nbd J\), and the annulus
\(\partial\nbd J\setminus\Int\nbd C_{w,t}(J)\) has both boundary curves of
slope \(wt\) on \(\partial\nbd C_{w,t}(J)\). (Several indexing
conventions for cables are in use; we fix the one above, in which the
first index is the winding number.)

\subsection{Fibered knots}\label{sec:fibered}

A knot \(K\) is \emph{fibered} if \(E_K\) is the total space of a fiber
bundle over \(S^1\) whose fibers are Seifert surfaces for \(K\), that is,
compact connected oriented surfaces \(F\subset E_K\) with
\(\partial F=\lambda\). The fiber then realises the minimal genus of a
Seifert surface for \(K\), so its genus is the \emph{genus} \(g=g(K)\)
of the knot, and \(g\geq1\) when \(K\) is non-trivial. Since \(F\) is a
compact connected orientable surface of genus \(g\) with one boundary
component, it deformation retracts onto a wedge of \(2g\) circles, and
\(\pi_1(F)\) is a free group \(F_{2g}\) of rank \(2g\). The long exact
homotopy sequence of the bundle gives a short exact sequence
\begin{equation}\label{eq:fibration}
 1\longrightarrow\pi_1(F)\longrightarrow\pi_1(E_K)
 \longrightarrow\pi_1(S^1)=\Z\longrightarrow1,
\end{equation}
because \(\pi_2(S^1)=0\) and \(F\) is connected. The meridian \(\mu\)
meets each fiber transversally in a single point, so the map to \(\Z\)
in \eqref{eq:fibration} sends \(\mu\) to a generator.

\subsection{Heegaard Floer homology and \texorpdfstring{\(L\)}{L}-spaces}\label{sec:Lspace}

Heegaard Floer homology \cite{OS2004a} assigns to a closed oriented
\(3\)-manifold \(Y\) a finite-dimensional vector space \(\HFhat(Y)\) over
\(\F_2\) (we use \(\F_2\) coefficients throughout), invariant up to
homeomorphism. If \(Y\) is a rational homology sphere, then
\(\dim_{\F_2}\HFhat(Y)\geq|H_1(Y;\Z)|\): the group splits as a direct
sum indexed by the spin\(^c\) structures of \(Y\), of which there are
\(|H^2(Y;\Z)|=|H_1(Y;\Z)|\), and each summand has Euler characteristic
\(\pm1\) \cite[Proposition~5.1]{OS2004b}. We call \(Y\) an
\emph{\(L\)-space} if equality holds. (The original definition
\cite[Definition~1.1]{OS2005} requires \(\HFhat(Y;\Z)\) to be a free
abelian group of rank \(|H_1(Y;\Z)|\); the \(\F_2\)-version above is
the one commonly used in the Heegaard Floer literature, and it is the one
used in this note.)
Lens spaces, including \(S^3\), are \(L\)-spaces
\cite[Proposition~3.1]{OS2004b}, which explains the name. The group
\(\HFhat(-Y)\) is the dual of \(\HFhat(Y)\)
\cite[Proposition~2.5]{OS2004b}, so \(Y\) is an \(L\)-space if and only
if \(-Y\) is.

A knot \(K\subset S^3\) is an \emph{\(L\)-space knot} if \(\Surg{r}{K}\)
is an \(L\)-space for some \(r\in\Q\); by (S2) such an \(r\) is
automatically non-zero. We say that \(K\) admits a \emph{positive}
\(L\)-space surgery if this holds for some \(r>0\). By (S4) and the
invariance of the \(L\)-space condition under orientation reversal,
\(K\) is an \(L\)-space knot if and only if \(K\) or its mirror image
admits a positive \(L\)-space surgery. The unknot is an \(L\)-space knot
by (S5), and so is every torus knot, since \((wt\pm1)\)-surgery on
\(T_{w,t}\) is a lens space \cite{Moser1971}.

\subsection{Group-theoretic tools}\label{sec:groups}

For a finitely generated group \(G\), let \(\rk(G)\) denote the minimal
number of generators of \(G\), and let \(G'=[G,G]\) denote its commutator
subgroup. If \(G\to Q\) is an epimorphism, then \(\rk(Q)\leq\rk(G)\).
Grushko's theorem \cite{Grushko1940} states that for finitely generated
groups \(A\) and \(B\),
\begin{equation}\label{eq:grushko}
 \rk(A*B)=\rk(A)+\rk(B).
\end{equation}
We use the following consequence of the Poincar\'e conjecture, proved
by Perelman (see \cite{MorganTian2007} for a detailed exposition).

\begin{lemma}\label{lem:perfect}
Let \(\Sigma\) be an integer homology sphere with \(\Sigma\not\cong S^3\).
Then \(\pi_1(\Sigma)\) is a non-trivial perfect group, and
\(\rk(\pi_1(\Sigma))\geq2\).
\end{lemma}

\begin{proof}
If \(\pi_1(\Sigma)\) were trivial, then \(\Sigma\) would be a simply
connected closed \(3\)-manifold, hence homeomorphic to \(S^3\) by the
Poincar\'e conjecture. Its abelianization is
\(H_1(\Sigma;\Z)=0\), so \(\pi_1(\Sigma)\) is perfect.
\review{It took a while to comprehend that ``Its'' is referring to $\pi_1(\Sigma)$. This should be written better.}

A non-trivial
cyclic group has non-trivial abelianization, so \(\pi_1(\Sigma)\) is not
cyclic, that is, \(\rk(\pi_1(\Sigma))\geq2\).
\end{proof}

A group \(G\) is \emph{normally generated} by a subset \(S\) if the
normal closure \(\langle\!\langle S\rangle\!\rangle\) is all of \(G\); the
\emph{weight} of \(G\) is the minimal size of such an \(S\). If \(G\) is
normally generated by one element and \(G\to Q\) is an epimorphism, then
\(Q\) is normally generated by the image of that element. By (S3), the
fundamental group of every Dehn surgery on a knot in \(S^3\) is normally
generated by the image of a meridian, that is, has weight at most one.
The group-theoretic input that limits the number of summands with
non-zero homology is the following theorem of Howie, previously known as
the Scott--Wiegold conjecture.

\begin{theorem}[Howie {\cite{Howie2002}}]\label{thm:howie}
Let \(a,b,c\geq2\). The free product
\((\Z/a\Z)*(\Z/b\Z)*(\Z/c\Z)\) is not normally generated by a single
element.
\end{theorem}

\section{The group-rank obstruction}\label{sec:rank}

In this section \(K\subset S^3\) is a fibered knot of genus \(g\), with
exterior \(E_K\) and fiber \(F\) as in Section~\ref{sec:fibered}. No
Floer-theoretic hypothesis is used here.

\begin{lemma}[Upper bound from the fiber]\label{lem:fiber}
Let \(K\subset S^3\) be a fibered knot of genus \(g\), and let
\(M=\Surg{r}{K}\) for any slope \(r\in\Q\cup\{\infty\}\). Then
\[
 \rk\bigl(\pi_1(M)'\bigr)\leq 2g.
\]
\end{lemma}

\begin{proof}
Put \(\Pi=\pi_1(E_K)\) and \(G=\pi_1(M)\). We first identify \(\Pi'\)
with the fiber group. Let \(\phi\colon\Pi\to\Z\) be the map in
\eqref{eq:fibration}, so that \(\ker\phi=\pi_1(F)\cong F_{2g}\). Since
\(\Z\) is abelian, \(\phi\) factors as \(\phi=\psi\circ\alpha\), where
\(\alpha\colon\Pi\to\Pi/\Pi'=H_1(E_K;\Z)\cong\Z\) is the abelianization
and \(\psi\colon\Z\to\Z\) is a homomorphism. Both \(\alpha\) and
\(\phi\) send the meridian \(\mu\) to a generator (by (S2) and
Section~\ref{sec:fibered}, respectively), so \(\psi\) is onto, hence
\(\psi=\pm\mathrm{id}\) and \(\ker\phi=\ker\alpha=\Pi'\). Thus
\(\Pi'=\pi_1(F)\cong F_{2g}\).

\review{The above should probably have made the list of standard facts.}

By (S3), Dehn filling induces an epimorphism \(f\colon\Pi\to G\). For
any epimorphism of groups, \(f(\Pi')=G'\): the inclusion
\(f(\Pi')\subseteq G'\) is clear, and conversely every commutator
\([x,y]\) in \(G\) equals \(f([\tilde x,\tilde y])\) for lifts
\(\tilde x,\tilde y\) of \(x,y\). Hence \(G'\) is a quotient of
\(F_{2g}\), and \(\rk(G')\leq2g\).
\end{proof}

The next lemma is the heart of the argument. The idea is that
\(\pi_1(M)'\) is the fundamental group of the universal abelian cover of
\(M\), and that in this cover a homology sphere summand of \(M\) is
replicated once for each sheet, each replica contributing at least two
generators.

\begin{lemma}[Lower bound from the universal abelian cover]\label{lem:cover}
Let \(M\) be a closed connected oriented \(3\)-manifold with
\(H_1(M;\Z)\cong\Z/n\Z\) for some \(n\geq1\), and suppose that
\[
 M\cong X\#\Sigma,
 \qquad H_1(\Sigma;\Z)=0,\qquad \Sigma\not\cong S^3,
\]
for some closed oriented \(3\)-manifold \(X\). Then
\[
 \rk\bigl(\pi_1(M)'\bigr)\geq 2n.
\]
In particular, if \(M=\Surg{p/q}{K}\) with \(p\neq0\) for a fibered knot
\(K\) of genus \(g\), then a summand \(\Sigma\) as above can occur only if
\begin{equation}\label{eq:numerator}
 |p|\leq g.
\end{equation}
\end{lemma}

\begin{proof}
Choose closed balls \(B_X\subset X\) and \(B_\Sigma\subset\Sigma\) and
write \(M=X^\circ\cup_S\Sigma^\circ\), where
\(X^\circ=X\setminus\Int B_X\) and \(\Sigma^\circ=\Sigma\setminus\Int B_\Sigma\)
are the summands with the interior of a ball removed and \(S\) is the
common boundary \(2\)-sphere. By \eqref{eq:vk},
\(H_1(M;\Z)\cong H_1(X;\Z)\oplus H_1(\Sigma;\Z)=H_1(X;\Z)\), so
\(H_1(X;\Z)\cong\Z/n\Z\) and the inclusion \(X^\circ\hookrightarrow M\)
induces an isomorphism \(H_1(X^\circ;\Z)\to H_1(M;\Z)\).

Let \(G=\pi_1(M)\) and let \(\rho\colon\widetilde M\to M\) be the
connected covering space corresponding to the normal subgroup \(G'\)
(the \emph{universal abelian cover} of \(M\))
\review{The phrase has already been used so it's not the time to define it.}

It is a regular covering
with deck group \(G/G'\cong H_1(M;\Z)\cong\Z/n\Z\),
so it has \(n\) sheets, and \(\pi_1(\widetilde M)\cong G'\). We use the
following standard fact about a regular covering with deck group \(Q\)
and a connected codimension-zero submanifold \(A\subseteq M\)
\cite[\S1.3]{HatcherAT}: if \(Q_A\) is the image of the composite
\(\pi_1(A)\to\pi_1(M)\to Q\), then \(\rho^{-1}(A)\) has exactly
\([Q:Q_A]\) components, and each component is the connected covering of
\(A\) associated with the kernel of that composite.

We apply this to the two pieces of \(M\). For \(A=X^\circ\), the composite
\(\pi_1(X^\circ)\to G\to G/G'\cong H_1(M;\Z)\) is the abelianization of
\(\pi_1(X^\circ)\) followed by the isomorphism
\(H_1(X^\circ;\Z)\to H_1(M;\Z)\) noted above.
\review{The above sentence is quite confusing. For a while I thought it was saying that the first map in the composition, the map $\pi_1(X^\circ)\to \pi_1(M)$, is the abelianization of $\pi_1(X^\circ)$. Ultimately it's correct, but not easy to parse.}
So it is onto, and its kernel
is \(\pi_1(X^\circ)'=\pi_1(X)'\). Hence \(\rho^{-1}(X^\circ)\) is
connected, and it is the covering of \(X^\circ\) associated with
\(\pi_1(X)'\). If \(\widetilde X\to X\) denotes the connected
\(n\)-sheeted covering associated with \(\pi_1(X)'\), this says that
\[
 \rho^{-1}(X^\circ)\cong\widetilde X^\circ:=\widetilde X\setminus
 (\text{the \(n\) lifts of }\Int B_X).
\]
For \(A=\Sigma^\circ\), the group \(\pi_1(\Sigma^\circ)=\pi_1(\Sigma)\) is
perfect by Lemma~\ref{lem:perfect},
\review{The group is perfect because $H_1$ is trivial, don't need lemma.}
so it maps trivially to the abelian
deck group \(G/G'\). Hence \(\rho^{-1}(\Sigma^\circ)\) consists of \(n\)
components, each mapped homeomorphically onto \(\Sigma^\circ\).

The \(n\) boundary spheres of \(\widetilde X^\circ\) are the lifts of
\(S\), and each of them is glued to one copy of \(\Sigma^\circ\). Since
\(\widetilde X^\circ\) is connected, this exhibits
\begin{equation}\label{eq:cover-sum}
 \widetilde M\cong
 \widetilde X\#\underbrace{\Sigma\#\cdots\#\Sigma}_{n\text{ copies}}.
\end{equation}
All fundamental groups in sight are finitely generated, being
fundamental groups of compact manifolds. By \eqref{eq:vk}, Grushko's
theorem \eqref{eq:grushko} and Lemma~\ref{lem:perfect},
\[
 \rk(G')=\rk\bigl(\pi_1(\widetilde M)\bigr)
 =\rk\bigl(\pi_1(\widetilde X)\bigr)+n\,\rk\bigl(\pi_1(\Sigma)\bigr)
 \geq 2n.
\]

Finally, let \(M=\Surg{p/q}{K}\) with \(p\neq0\) and \(K\) fibered of
genus \(g\). By (S2), \(H_1(M;\Z)\cong\Z/|p|\Z\), so the first part applies
with \(n=|p|\) and gives \(\rk(G')\geq2|p|\), while
Lemma~\ref{lem:fiber} gives \(\rk(G')\leq2g\). Hence \(|p|\leq g\).
\end{proof}

Howie's theorem (Theorem~\ref{thm:howie}) limits the summands with
non-zero homology; this is what produces the homology sphere summand to
which Lemma~\ref{lem:cover} applies.

\begin{lemma}[Consequence of Howie's theorem {\cite{Howie2002}}]\label{lem:howie}
Let \(K\subset S^3\) be any knot and let \(p/q\) be a slope with
\(p\neq0\). Then at most two non-trivial prime summands of
\(\Surg{p/q}{K}\) have non-zero first integral homology.
\end{lemma}

\begin{proof}
Set \(M=\Surg{p/q}{K}\) and let \(M_1\#\cdots\#M_k\) be its decomposition
into non-trivial prime summands. By \eqref{eq:vk} and (S2),
\[
 \bigoplus_{i=1}^k H_1(M_i;\Z)\cong H_1(M;\Z)\cong\Z/|p|\Z,
\]
so each \(H_1(M_i;\Z)\) is a direct summand of a finite cyclic group and
hence is finite cyclic. Suppose that three of them are non-zero, of
orders \(a,b,c\geq2\). By \eqref{eq:vk},
\(\pi_1(M)\cong\pi_1(M_1)*\cdots*\pi_1(M_k)\). Abelianizing those three
free factors and killing all the others defines an epimorphism
\[
 \pi_1(M)\twoheadrightarrow(\Z/a\Z)*(\Z/b\Z)*(\Z/c\Z).
\]
By (S3), \(\pi_1(M)\) is normally generated by one element, and this
property passes to the quotient (Section~\ref{sec:groups}). This
contradicts Theorem~\ref{thm:howie}.
\end{proof}

\begin{proposition}\label{prop:large}
Let \(K\subset S^3\) be a fibered knot of genus \(g\geq1\). If
\(p\neq0\) and \(|p|\geq2g-1\), then \(\Surg{p/q}{K}\) has at most two
non-trivial prime summands.
\end{proposition}

\begin{proof}
Suppose, for a contradiction, that \(M=\Surg{p/q}{K}\cong M_1\#\cdots\#M_k\)
with \(k\geq3\) and each \(M_i\) a prime manifold other than \(S^3\). By
Lemma~\ref{lem:howie}, at most two of the \(M_i\) have non-zero first
homology, so after renumbering \(H_1(M_k;\Z)=0\); thus \(M_k\) is an
integer homology sphere (Section~\ref{sec:prime}) with
\(M_k\not\cong S^3\). Apply Lemma~\ref{lem:cover} with
\(X=M_1\#\cdots\#M_{k-1}\) and \(\Sigma=M_k\): since \(p\neq0\), it gives
\(|p|\leq g\), and together with the hypothesis,
\[
 2g-1\leq|p|\leq g .
\]
This forces \(g\leq1\), so \(g=1\) and \(|p|=1\).

It remains to rule out this endpoint. If \(|p|=1\), then
\(H_1(M;\Z)=0\) by (S2), so by \eqref{eq:vk} every \(M_i\) is an integer
homology sphere other than \(S^3\). Put \(G=\pi_1(M)\). Then \(G\) is
perfect, so \(G'=G\), and \eqref{eq:vk}, Grushko's theorem
\eqref{eq:grushko} and Lemma~\ref{lem:perfect} give
\[
 \rk(G')=\rk(G)=\sum_{i=1}^k\rk\bigl(\pi_1(M_i)\bigr)\geq2k\geq6.
\]
But Lemma~\ref{lem:fiber} gives \(\rk(G')\leq2g=2\), a contradiction.
\end{proof}

Proposition~\ref{prop:large} is Theorem~\ref{thm:fibered}
\review{Then why two different names?}
Note that it
makes no assumption on the knot beyond fiberedness. The proof actually
excludes a homology sphere summand whenever \(|p|>g\), by
\eqref{eq:numerator}. The threshold \(|p|\geq 2g-1\) in the statement is
the one that the Hom--Lidman--Zufelt theorem will supply for hyperbolic
\(L\)-space knots; for \(g=1\) it admits the value \(|p|=1\), which is
why the endpoint required the separate argument above.
\review{Garbled sentence: what's the antecedent of ``it''?}

\begin{remark}\label{rem:algebraic}
Lemma~\ref{lem:cover} has a purely group-theoretic core: if \(G=A*H\)
with \(H\) perfect and \(A/A'\cong\Z/n\Z\), then
\(G'\cong A'*H*\cdots*H\) with \(n\) copies of \(H\). This can also be
derived from the Kurosh subgroup theorem \cite{Kurosh1934}; the covering
space argument above makes the multiplicity \(n\) visible directly.
\end{remark}

\begin{example}\label{ex:torus}
Lemma~\ref{lem:fiber} is sharp: for a torus knot, the cabling surgery
already exhausts the bound \(\rk(G')\leq2g\) with a free group, leaving
no room for a homology sphere summand. Let \(T_{a,b}\) be a torus knot
with \(a,b\geq2\) coprime, so that \(2g(T_{a,b})=(a-1)(b-1)\), and recall
that \(\Surg{ab}{T_{a,b}}\cong L(a,b)\#L(b,a)\) by Moser's theorem
\cite{Moser1971}. The fundamental group is \(G=(\Z/a\Z)*(\Z/b\Z)\), and
the kernel \(G'\) of its abelianization onto \(\Z/ab\Z\) is a free group
of rank \((a-1)(b-1)=2g\): indeed, in the \(ab\)-fold universal abelian
cover the punctured \(L(a,b)\) lifts to \(b\) copies of \(S^3\) with
\(a\) balls removed, the punctured \(L(b,a)\) lifts to \(a\) copies of
\(S^3\) with \(b\) balls removed, and these are glued along the \(ab\)
lifts of the separating sphere; the resulting manifold is a connected
sum of \(ab-(a+b)+1=(a-1)(b-1)\) copies of \(S^1\times S^2\). This is
consistent with \eqref{eq:numerator}, since \(ab>g\).
\end{example}

\section{Inputs from Heegaard Floer theory and Dehn surgery}\label{sec:inputs}

We now state the published theorems used in the proof of
Theorem~\ref{thm:main}. None of them is reproved here.

\begin{theorem}[Ni {\cite[Corollary~1.3]{Ni2007}}]\label{thm:ni}
Every \(L\)-space knot in \(S^3\) is fibered.
\end{theorem}

Ni's corollary explicitly allows a rational \(L\)-space surgery slope,
and its proof uses the structure theorem of Ozsv\'ath and Szab\'o
\cite{OS2005} for the knot Floer homology (the knot-level refinement of
\(\HFhat\)) of \(L\)-space knots; see also
Ni's erratum \cite{NiCorrigendum}, which repairs part of the proof of
the fiberedness detection theorem without affecting the statement, and
Ghiggini \cite{Ghiggini2008} for the genus-one case.

\begin{theorem}[Hom--Lidman--Zufelt {\cite[Theorem~1.3]{HLZ2015}}]\label{thm:hlz}
Let \(K\subset S^3\) be a hyperbolic knot that admits a positive
\(L\)-space surgery. If \(\Surg{r}{K}\) is reducible, then
\(r=2g(K)-1\).
\end{theorem}

Two features of this statement matter for us. First, the reducible
surgery itself is not assumed to be an \(L\)-space; the theorem
constrains every reducible slope of \(K\). Second, the hyperbolicity
hypothesis cannot be dropped: the right-handed trefoil \(T_{2,3}\) has
genus \(1\) and admits a positive \(L\)-space surgery (its \(5\)- and
\(7\)-surgeries are lens spaces \cite{Moser1971}), but its reducible
surgery \(\Surg{6}{T_{2,3}}\cong L(2,3)\#L(3,2)\) has slope
\(6\neq 2g-1\).
This is why the proof of Theorem~\ref{thm:main} treats non-hyperbolic
knots separately, using the following known cases of the Cabling
Conjecture.

\begin{theorem}[Moser {\cite{Moser1971}}, Scharlemann {\cite[Corollaries~4.4--4.5]{Scharlemann1990}}]\label{thm:cabling}
Let \(K\subset S^3\) be a non-trivial knot that is not hyperbolic, that
is, a torus knot or a satellite knot. If \(\Surg{r}{K}\) is reducible,
then \(K=C_{w,t}(J)\) is a cable, with \(w\geq2\) and \(\gcd(w,t)=1\), of
some knot \(J\), and \(r=wt\) is the cabling slope.
\end{theorem}

For torus knots this is Moser's classification of surgeries on
\(T_{w,t}=C_{w,t}(U)\): only the slope \(wt\) yields a reducible
manifold. For satellite knots it is Scharlemann's theorem; Corollary~4.4
of \cite{Scharlemann1990} treats a knot in a solid torus that is not
contained in a ball, and Corollary~4.5 applies it to satellite knots in
\(S^3\). Since composite knots are satellite knots
(Section~\ref{sec:types}), they are covered. The reduction of the
Cabling Conjecture to hyperbolic knots along these lines is standard and
is recorded, for instance, in the introduction of \cite{HLZ2015}.

\begin{theorem}[Gordon {\cite[Corollary~7.3]{Gordon1983}}]\label{thm:gordon}
Let \(J\subset S^3\) be a knot and let \(w\geq2\), \(t\) be coprime
integers. Then, up to the orientation of the lens space,
\begin{equation}\label{eq:cable}
 \Surg{wt}{C_{w,t}(J)}\cong L(w,t)\#\Surg{t/w}{J}.
\end{equation}
\end{theorem}

Formula \eqref{eq:cable} is consistent with (S2): both sides have first
homology of order \(|wt|=w\cdot|t|\), since \(\gcd(w,t)=1\). Geometrically,
the cabling annulus together with two meridian discs of the filling
solid torus forms a separating sphere exhibiting the decomposition
\eqref{eq:cable}: on one side it cuts off a punctured \(\Surg{t/w}{J}\),
and on the other a punctured lens space of order \(w\), obtained from
the solid torus \(\nbd J\) by attaching a \(2\)-handle along a curve of
winding number \(w\). This sphere is a reducing sphere when both factors
are non-trivial, but not always: for \(J\) the unknot and \(|t|=1\) the
cable \(C_{w,t}(J)\) is itself the unknot and \(\Surg{t/w}{J}\cong S^3\),
so the sphere bounds a ball.

\begin{theorem}[Gordon--Luecke {\cite{GordonLuecke1987}}]\label{thm:gl}
If \(J\subset S^3\) is a non-trivial knot and \(r\in\Q\) is not an
integer, then \(\Surg{r}{J}\) is irreducible.
\end{theorem}

\section{Proof of Theorem~\ref{thm:main}}\label{sec:proof}

Let \(K\subset S^3\) be an \(L\)-space knot and \(r\in\Q\cup\{\infty\}\).
We must show that \(\Surg{r}{K}\) has at most two non-trivial prime
summands.

\paragraph{Step 1: reduction to a positive \(L\)-space surgery.}
By the last paragraph of Section~\ref{sec:Lspace}, either \(K\) or its
mirror image \(\overline K\) admits a positive \(L\)-space surgery. By (S4),
\(\Surg{-r}{\overline K}\cong-\Surg{r}{K}\), and by
Section~\ref{sec:prime} a manifold and its orientation reversal have the
same number of non-trivial prime summands. Mirroring also preserves the
genus and the type (torus, satellite or hyperbolic) of a knot. Replacing
\((K,r)\) by \((\overline K,-r)\) if necessary, we may therefore assume
that \(K\) admits a positive \(L\)-space surgery.

\review{A little heavy on the details, e.g. ``mirroring preserves knot type'' is too pedantic to mention. This is more of a preliminary remark than a step.}

\paragraph{Step 2: the unknot.}
If \(K\) is the unknot, then by (S5) \(\Surg{r}{K}\) is a lens space,
\(S^3\) or \(S^1\times S^2\). Lens spaces are irreducible and
\(S^1\times S^2\) is prime (Section~\ref{sec:prime}), so there is at most
one non-trivial prime summand. Henceforth \(K\) is non-trivial; by
Theorem~\ref{thm:ni} it is fibered, of genus \(g\geq1\).

\paragraph{Step 3: the slopes \(\infty\) and \(0\).}
By (S1), \(\Surg{\infty}{K}=S^3\) has no non-trivial prime summand. For
the slope \(0\), the filling solid torus has meridian \(\lambda=\partial F\),
where \(F\) is a fiber; the fibration of \(E_K\) restricts on
\(\partial E_K\) to the product fibration by longitudes, so the solid
torus is a union of discs capping off the fibers, and \(\Surg{0}{K}\) is a
bundle over \(S^1\) with fiber the closed surface \(\widehat F=F\cup D^2\)
of genus \(g\geq1\). The universal cover of \(\widehat F\) is \(\R^2\),
so the pull-back of the bundle to the universal cover \(\R\) of \(S^1\)
is \(\widehat F\times\R\), and the universal cover of \(\Surg{0}{K}\) is
\(\R^3\). Hence \(\Surg{0}{K}\) is irreducible (Section~\ref{sec:prime};
\cite[Proposition~1.6]{Hatcher3M}), so it is prime and has exactly one
non-trivial prime summand (it is not \(S^3\), having \(H_1\cong\Z\)).

\paragraph{Step 4: finite non-zero slopes.}
Let \(r=p/q\) with \(p\neq0\). If \(\Surg{r}{K}\) is prime, it has at most
one non-trivial prime summand. Otherwise it is reducible, and we
distinguish the type of \(K\) (Section~\ref{sec:types}).

\emph{Case 4a: \(K\) is hyperbolic.} By Step~1, \(K\) admits a positive
\(L\)-space surgery, so Theorem~\ref{thm:hlz} gives \(r=2g-1\), that is,
\(p=2g-1\geq1\) and \(q=1\). Since \(K\) is fibered of genus \(g\geq1\),
Proposition~\ref{prop:large} applies and shows that \(\Surg{r}{K}\) has at
most two non-trivial prime summands.

\emph{Case 4b: \(K\) is a torus knot or a satellite knot.} By
Theorem~\ref{thm:cabling}, \(K=C_{w,t}(J)\) with \(w\geq2\),
\(\gcd(w,t)=1\), and \(r=wt\). By Theorem~\ref{thm:gordon},
\(\Surg{r}{K}\cong L(w,t)\#\Surg{t/w}{J}\). The first factor is a lens
space with \(|H_1|=w\geq2\), hence a single non-trivial prime summand.
The slope \(t/w\) is not an integer, because \(w\geq2\) and
\(\gcd(w,t)=1\). If \(J\) is non-trivial, \(\Surg{t/w}{J}\) is irreducible
by Theorem~\ref{thm:gl}, hence prime, and contributes at most one
non-trivial prime summand. If \(J\) is the unknot, then \(t\neq0\) (as
\(\gcd(w,0)=w\geq2\)), so by (S5) \(\Surg{t/w}{J}\) is a lens space or
\(S^3\), and again contributes at most one. By uniqueness of the prime
decomposition, \(\Surg{r}{K}\) has at most two non-trivial prime
summands.

This completes the proof of Theorem~\ref{thm:main}. \qed

\section{Comments}\label{sec:comments}

\begin{remark}[Where the \(L\)-space hypothesis is used]
The proof uses the \(L\)-space hypothesis only through
Theorems~\ref{thm:ni} and~\ref{thm:hlz}, in Steps~2 and~4a (Step~1
merely arranges the sign of the \(L\)-space slope that
Theorem~\ref{thm:hlz} requires). Consequently the same argument proves
the following statement, of which Theorem~\ref{thm:main} is a special
case: if \(K\) is a non-trivial fibered knot that is either
non-hyperbolic, or hyperbolic with every reducible slope \(p/q\)
satisfying \(|p|\geq2g(K)-1\), then every surgery on \(K\) has at most
two non-trivial prime summands. (Fiberedness is used in Step~3 for the
slope \(0\) and in Step~4a through Proposition~\ref{prop:large}.)
\end{remark}

\begin{remark}[What is not proved]\label{rem:notproved}
The argument does not establish the Cabling Conjecture for hyperbolic
\(L\)-space knots. In Case~4a it handles a hypothetical reducible surgery
at slope \(2g-1\) directly, by excluding the integer homology sphere
summand that a third summand would force, without attempting to identify
that summand by Heegaard Floer methods. It does not exclude a reducible
surgery \(\Surg{2g-1}{K}\) on a hyperbolic \(L\)-space knot with exactly
two non-trivial prime summands, both with non-zero first homology.
For an arbitrary knot in \(S^3\), a surgery has at most three
non-trivial prime summands, and three summands would have to be two lens
spaces of coprime orders and an integer homology sphere; this combines
results of Sayari, Valdez-S\'anchez and Howie, see the introduction of
\cite{Howie2010}. 
\review{Ok but cite the original sourcess too.}
Proposition~\ref{prop:large} excludes the last
configuration in the range \(|p|\geq2g-1\) for fibered knots without
using this classification.
\end{remark}

\begin{remark}[The ambient manifold]
The hypothesis that \(K\) lies in \(S^3\) is used throughout: in (S2) and
(S3), in Ni's and Hom--Lidman--Zufelt's theorems, and in the Cabling
Conjecture. For knots in other \(L\)-spaces, a surgery can have three
non-trivial prime summands: a regular fiber of the Seifert fibration of
the Poincar\'e homology sphere (itself an \(L\)-space) has a surgery
homeomorphic to \(L(2,1)\#L(3,2)\#L(5,4)\) \cite{Caudell2021}.
\end{remark}

\begin{thebibliography}{99}

\bibitem{BodishDeYeso}
H.~Bodish and R.~DeYeso~III,
\emph{Obstructing reducible surgeries: slice genus and thickness bounds},
arXiv:2209.01672 (2022).

\bibitem{Caudell2021}
J.~Caudell,
\emph{Three lens space summands from the Poincar\'e homology sphere},
arXiv:2101.01256 (2021).

\bibitem{GAS1986}
F.~Gonz\'alez-Acu\~na and H.~Short,
\emph{Knot surgery and primeness},
Math.\ Proc.\ Cambridge Philos.\ Soc.\ \textbf{99} (1986), no.~1, 89--102.

\bibitem{Ghiggini2008}
P.~Ghiggini,
\emph{Knot Floer homology detects genus-one fibred knots},
Amer.\ J.\ Math.\ \textbf{130} (2008), no.~5, 1151--1169.

\bibitem{Gordon1983}
C.~McA.~Gordon,
\emph{Dehn surgery and satellite knots},
Trans.\ Amer.\ Math.\ Soc.\ \textbf{275} (1983), no.~2, 687--708.

\bibitem{GordonLuecke1987}
C.~McA.~Gordon and J.~Luecke,
\emph{Only integral Dehn surgeries can yield reducible manifolds},
Math.\ Proc.\ Cambridge Philos.\ Soc.\ \textbf{102} (1987), no.~1, 97--101.

\bibitem{Grushko1940}
I.~A.~Grushko,
\emph{On the bases of a free product of groups},
Mat.\ Sbornik (N.S.) \textbf{8(50)} (1940), no.~1, 169--182 (Russian).

\bibitem{HatcherAT}
A.~Hatcher,
\emph{Algebraic topology},
Cambridge University Press, 2002.

\bibitem{Hatcher3M}
A.~Hatcher,
\emph{Notes on basic \(3\)-manifold topology},
unpublished notes, revised 2000.

\bibitem{HLZ2015}
J.~Hom, T.~Lidman and N.~Zufelt,
\emph{Reducible surgeries and Heegaard Floer homology},
Math.\ Res.\ Lett.\ \textbf{22} (2015), no.~3, 763--788; arXiv:1307.5317.

\bibitem{Howie2002}
J.~Howie,
\emph{A proof of the Scott--Wiegold conjecture on free products of
cyclic groups},
J.\ Pure Appl.\ Algebra \textbf{173} (2002), no.~2, 167--176.

\bibitem{Howie2010}
J.~Howie,
\emph{Can Dehn surgery yield three connected summands?},
Groups Geom.\ Dyn.\ \textbf{4} (2010), no.~4, 785--797; arXiv:0908.2697.

\bibitem{Kurosh1934}
A.~Kurosch,
\emph{Die Untergruppen der freien Produkte von beliebigen Gruppen},
Math.\ Ann.\ \textbf{109} (1934), 647--660.

\bibitem{MorganTian2007}
J.~Morgan and G.~Tian,
\emph{Ricci flow and the Poincar\'e conjecture},
Clay Mathematics Monographs \textbf{3}, Amer.\ Math.\ Soc., 2007.

\bibitem{Moser1971}
L.~Moser,
\emph{Elementary surgery along a torus knot},
Pacific J.\ Math.\ \textbf{38} (1971), no.~3, 737--745.

\bibitem{Ni2007}
Y.~Ni,
\emph{Knot Floer homology detects fibred knots},
Invent.\ Math.\ \textbf{170} (2007), no.~3, 577--608.

\bibitem{NiCorrigendum}
Y.~Ni,
\emph{Erratum: Knot Floer homology detects fibred knots},
Invent.\ Math.\ \textbf{177} (2009), no.~1, 235--238; arXiv:0808.0940.

\bibitem{OS2004a}
P.~Ozsv\'ath and Z.~Szab\'o,
\emph{Holomorphic disks and topological invariants for closed
three-manifolds},
Ann.\ of Math.\ (2) \textbf{159} (2004), no.~3, 1027--1158.

\bibitem{OS2004b}
P.~Ozsv\'ath and Z.~Szab\'o,
\emph{Holomorphic disks and three-manifold invariants: properties and
applications},
Ann.\ of Math.\ (2) \textbf{159} (2004), no.~3, 1159--1245.

\bibitem{OS2005}
P.~Ozsv\'ath and Z.~Szab\'o,
\emph{On knot Floer homology and lens space surgeries},
Topology \textbf{44} (2005), no.~6, 1281--1300.

\bibitem{Scharlemann1990}
M.~Scharlemann,
\emph{Producing reducible \(3\)-manifolds by surgery on a knot},
Topology \textbf{29} (1990), no.~4, 481--500.

\bibitem{Thurston1982}
W.~P.~Thurston,
\emph{Three dimensional manifolds, Kleinian groups and hyperbolic
geometry},
Bull.\ Amer.\ Math.\ Soc.\ (N.S.) \textbf{6} (1982), no.~3, 357--381.

\end{thebibliography}

\end{document}
